\documentclass[11pt]{article}

\usepackage{amsmath,amssymb}
\usepackage{xurl}
\usepackage[colorlinks=true,linkcolor=blue,citecolor=blue,urlcolor=blue]{hyperref}

% Shared commands for notes in docs/paper-gaps/.

\DeclareUnicodeCharacter{2080}{\ifmmode{}_{0}\else\textsubscript{0}\fi}
\DeclareUnicodeCharacter{2081}{\ifmmode{}_{1}\else\textsubscript{1}\fi}
\DeclareUnicodeCharacter{2082}{\ifmmode{}_{2}\else\textsubscript{2}\fi}
\DeclareUnicodeCharacter{2099}{\ifmmode{}_{n}\else\textsubscript{n}\fi}

\newcommand{\Lean}{\textsc{Lean}}
\ExplSyntaxOn
\NewDocumentCommand{\leanid}{m}{
  \ifmmode
    \textsf{\detokenize{#1}}
  \else
    \group_begin:
    \urlstyle{sf}
    \tl_set:Nn \l_tmpa_tl {#1}
    \tl_replace_all:Nnn \l_tmpa_tl {\_} {_}
    \exp_args:NV \path \l_tmpa_tl
    \group_end:
  \fi
}
\ExplSyntaxOff
\providecommand{\tr}{\operatorname{tr}}
\newcommand{\tnleanIssueBase}{https://github.com/LionSR/TNLean/issues/}
\newcommand{\tnleanPrBase}{https://github.com/LionSR/TNLean/pull/}
\newcommand{\tnleanDocBase}{https://sirui-lu.com/TNLean/docs/}
\newcommand{\ghissue}[1]{\href{\tnleanIssueBase#1}{GitHub issue \##1}}
\newcommand{\ghpr}[1]{\href{\tnleanPrBase#1}{PR \##1}}
\newcommand{\doclink}[1]{\href{\tnleanDocBase#1.html}{\path{#1}}}

% Dirac notation and the identity operator, as in blueprint/src/macros/common.tex.
% \providecommand keeps note-local definitions authoritative.
\usepackage{dsfont}
\providecommand{\ket}[1]{|#1\rangle}
\providecommand{\bra}[1]{\langle #1|}
\providecommand{\braket}[2]{\langle #1 | #2 \rangle}
\providecommand{\ketbra}[2]{|#1\rangle\!\langle #2|}
\providecommand{\Id}{\mathds{1}}
\providecommand{\one}{\mathds{1}}
\providecommand{\C}{\mathbb{C}}
\providecommand{\MN}[1]{M_{#1}(\C)}
\providecommand{\MD}{M_D(\C)}

% External referencing for paper sources.  The build helper writes sanitized
% auxiliary files under build/paper-aux/ and excludes duplicated source labels.
\usepackage{xr}

% Import a generated paper auxiliary file with a caller-supplied label prefix.
% The candidate paths support builds from docs/paper-gaps/, the repository root,
% and build/paper-aux/ itself.
\newcommand{\importpaperaux}[2]{%
  \IfFileExists{../../build/paper-aux/#2.aux}{%
    \externaldocument[#1]{../../build/paper-aux/#2}%
  }{%
    \IfFileExists{build/paper-aux/#2.aux}{%
      \externaldocument[#1]{build/paper-aux/#2}%
    }{%
      \IfFileExists{#2.aux}{%
        \externaldocument[#1]{#2}%
      }{}%
    }%
  }%
}

% General external reference macros
\newcommand{\paperref}[2]{\ref{#1-#2}}
\newcommand{\papereqref}[2]{(\ref{#1-#2})}

% CPSV16 source (1606.00608 / MPDO-22-12-17-2.tex) external document and shortcuts
\importpaperaux{cpsv16-}{1606.00608/MPDO-22-12-17-2-tnlean}
\newcommand{\cpsvref}[1]{\paperref{cpsv16}{#1}}
\newcommand{\cpsveqref}[1]{\papereqref{cpsv16}{#1}}

% Verdict marker: \gapnote{<kind>}{<status>}. Kind is the policy
% classification (clarification, local-correction, scope-restriction,
% unfaithful, false-source, open-gap); status is open, wip (elimination
% underway), resolved, or historical. Machine-read by the site index;
% prints a verdict line.
\newcommand{\gapnote}[2]{\par\noindent\textbf{Verdict:} #1 (#2).\par}


\title{Scope Note: The AKLT On-Site $\mathrm{SO}(3)$ Symmetry Is Formalized;
Its Cohomological Classification Is Not}
\author{}
\date{2026-06-04}

\begin{document}
\maketitle
\gapnote{scope-restriction}{open}

\section{The assertion}

The matrix-product-state review of Cirac, P\'erez-Garc\'ia, Schuch, and
Verstraete states that the Affleck--Kennedy--Lieb--Tasaki (AKLT) MPS lies in a
nontrivial symmetry-protected topological phase
\cite{Cirac2021Matrix}. It says that this phase ``can be protected by multiple
distinct physical symmetries: on-site $\mathrm{SO}(3)$ (with virtual symmetry
given by a spin-$\tfrac{1}{2}$ representation), on-site
$\mathbb{Z}_2\times\mathbb{Z}_2$ (the smallest subgroup of $\mathrm{SO}(3)$
still exhibiting a non-trivial 2-cocycle), time-reversal, or reflection
symmetry.''\footnote{See Ref.~\cite{Cirac2021Matrix}, around line~1159 of the
source.}

In each case, the physical symmetry acts on a spin-$1$ site, while the virtual
degrees of freedom carry a projective representation. The review identifies
the corresponding cohomology class with the order-two generator of the second
cohomology group of the symmetry group with $\mathrm{U}(1)$ coefficients
\cite{Cirac2021Matrix}.

The strongest statement in this list concerns the on-site $\mathrm{SO}(3)$
action. The physical index carries the spin-$1$ representation, while the
virtual index carries the projective spin-$\tfrac{1}{2}$ representation. Its
class is asserted to generate
\begin{align}
    H^2\!\left(\mathrm{SO}(3),\mathrm{U}(1)\right)
        &\cong \mathbb{Z}_2.
    \label{eq:rmp_aklt_continuous_rotation_gap_so3_cohomology}
\end{align}

\section{The point at issue}

The on-site $\mathrm{SO}(3)$ symmetry and the cohomological classification of
the protected phase are distinct mathematical statements. The first is
formalized; the second requires infrastructure that is not currently
available.

\paragraph{On-site symmetry.}

Let the Cartesian AKLT tensor be $A=(A^i)_{i=1}^3$, with
$A^i=\sigma^i$ given by the Pauli matrices. For each
$R\in\mathrm{SO}(3)$, applying $R$ to the physical index gives a tensor gauge
equivalent to $A$. The gauge is an $\mathrm{SU}(2)$ preimage $U$ of $R$ under
the adjoint double cover.\footnote{Formal declaration:
\leanid{aklt\_isOnSiteSymmetric\_SO3} in
\doclink{TNLean/MPS/Examples/AKLTRotation}.}

The substantive group-theoretic input is the surjectivity of the adjoint map
from $\mathrm{SU}(2)$ to $\mathrm{SO}(3)$. Conjugation by
$U\in\mathrm{SU}(2)$ rotates the Pauli vector:
\begin{align}
    U\sigma^iU^\dagger
        &= \sum_{j=1}^3 R(U)_{ij}\sigma^j.
    \label{eq:rmp_aklt_continuous_rotation_gap_pauli_adjoint_action}
\end{align}
Every rotation has a preimage:
\begin{align}
    \forall R\in\mathrm{SO}(3)\ \exists U\in\mathrm{SU}(2),
    \qquad
    R(U)
        &=R.
    \label{eq:rmp_aklt_continuous_rotation_gap_spin_half_cover_surjective}
\end{align}
The proof uses an explicit $ZXZ$ Euler-angle construction. Each coordinate
rotation is the image of an explicit $\mathrm{SU}(2)$ generator, and every
rotation factors into coordinate rotations.\footnote{Formal declaration:
\leanid{spinHalfCover\_surjective\_onto\_SO3}.}

The physical spin-$1$ representation is the defining action of
$\mathrm{SO}(3)$ on the three Cartesian components. The map
$U\mapsto R(U)$ is a monoid homomorphism, so its image is a subgroup. Showing
that the image contains the Euler generators proves
(\ref{eq:rmp_aklt_continuous_rotation_gap_spin_half_cover_surjective}) without
using Lie-group topology, a Clifford spin group, or a quaternion
identification; the library supplies none of those alternatives.

\paragraph{Cohomological classification.}

The assertion that the virtual projective representation realizes the
nontrivial element of
(\ref{eq:rmp_aklt_continuous_rotation_gap_so3_cohomology}) is separate from
the symmetry theorem. It requires continuous group cohomology in a Borel or
Lie-group formulation, or equivalently the Schur multiplier of
$\mathrm{SO}(3)$.

The group-cohomology library currently treats $\mathrm{SO}(3)$ as a discrete
group. That algebraic theory computes a different object from the continuous
cohomology needed here. The library also lacks a Schur-multiplier computation
and a projective-representation interface. Its developing
continuous-cohomology framework has no relevant computations and has not been
reconciled with the discrete theory. Formalizing this classification is
therefore a substantial independent development.

\section{The formalized statement}

The formalization proves the full on-site $\mathrm{SO}(3)$ symmetry of the
Cartesian AKLT tensor. The physical site carries the defining spin-$1$
representation, and the bond index is conjugated by a spin-$\tfrac{1}{2}$
$\mathrm{SU}(2)$ preimage of each rotation. This proves the symmetry statement
in Ref.~\cite{Cirac2021Matrix}, including the required cover surjectivity.

The finite subgroup $\mathbb{Z}_2\times\mathbb{Z}_2$, which the review
identifies as the smallest subgroup of $\mathrm{SO}(3)$ with a nontrivial
$2$-cocycle, is formalized separately on the $\ket{m}$-basis AKLT tensor
\cite{Cirac2021Matrix}. Its two gauges are $i\sigma_y$ and $\sigma_z$, and
they anticommute.\footnote{Formal declaration:
\leanid{aklt\_isOnSiteSymmetric\_Z2Z2} in
\doclink{TNLean/MPS/Examples/AKLT}.} Explicitly,
\begin{align}
    (i\sigma_y)\sigma_z
        &=-\sigma_z(i\sigma_y).
    \label{eq:rmp_aklt_continuous_rotation_gap_z2z2_anticommutator}
\end{align}

In both symmetry realizations, the bond action is projective. For
$\mathbb{Z}_2\times\mathbb{Z}_2$, the sign in
(\ref{eq:rmp_aklt_continuous_rotation_gap_z2z2_anticommutator}) is the explicit
obstruction. For the full $\mathrm{SO}(3)$ action, commuting rotations may
have $\mathrm{SU}(2)$ preimages that do not commute; the sign obstruction is
the nontriviality of the double cover. This obstruction is the physical
content associated with the symmetry-protected phase, but identifying its
continuous cohomology class remains outside the present formalization.

\section{Conclusion}

The on-site $\mathrm{SO}(3)$ symmetry of the AKLT tensor asserted in
Ref.~\cite{Cirac2021Matrix} is formalized, together with the surjectivity of
the spin-$\tfrac{1}{2}$ cover on which the proof depends. The unformalized
claim is the classification of the resulting projective phase as the
order-two generator of
$H^2(\mathrm{SO}(3),\mathrm{U}(1))\cong\mathbb{Z}_2$. Proving that statement
requires continuous group cohomology for $\mathrm{SO}(3)$, or an equivalent
Schur-multiplier theory, neither of which the current library provides.

The blueprint therefore links the symmetry and cover-surjectivity statements
to their formal declarations. It retains the cohomological classification as
an untagged remark.

\bibliographystyle{alpha}
\bibliography{references}

\end{document}
